% !TeX root = ../../Book2.tex
% 纲要标题：### 10.1 双线性型
% 纲要位置：计划纲要.md，第 1096 行。

\section{双线性型}
\label{b2:sec:1001}

线性映射接受一个向量，双线性型则同时比较两个向量。选定基以后，两者都能写成矩阵，但换基时的规则不同。本节先把配对两端的空间分开；只有在说明对称性或交替性以后，才把左右正交条件合并。

% 纲要内容（仅作写作计划，尚非教材正文）：
%
% * 配对、非退化性及正交补
% * 对称型与交替型
% * 合同分类与相似分类的区别
%

\subsection{配对、非退化性与正交补}
\label{b2:subsec:1001:01}

\begin{definition}\label{b2:def:bilinear-pairing}
设 \(k\) 为域，\(U,V\) 为 \(k\) 向量空间。如果映射 \(b:U\times V\to k\) 在固定任意一个变量后对另一个变量都是 \(k\) 线性的，就称 \(b\) 为\term[shuangxianxing-peidui]{双线性配对}[bilinear pairing]。它的左根空间与右根空间分别为
\[
\{\,u\in U\mid b(u,v)=0\text{ 对所有 }v\in V\,\},\qquad
\{\,v\in V\mid b(u,v)=0\text{ 对所有 }u\in U\,\}.
\]
两者均为零时，称配对\term[peidui-feituihua]{非退化}[nondegenerate]。当 \(U=V\) 时，称 \(b\) 为 \(V\) 上的\term[shuangxianxing-xing]{双线性型}[bilinear form]。
\end{definition}

以下空间均有限维。配对给出两个映射
\[
L_b:U\longrightarrow V^*,\quad u\longmapsto b(u,-),
\qquad
R_b:V\longrightarrow U^*,\quad v\longmapsto b(-,v).
\]
它们的核就是上述两个根空间。对这样的有限维 \(k\) 向量空间上的配对 \(b:U\times V\to k\)，若 \(L_b:U\to V^*\) 是同构，就称 \(b\) 为\term[wanmei-peidui]{完美配对}[perfect pairing]。

\begin{proposition}\label{b2:prop:finite-perfect-pairing}
设 \(k\) 为域，\(U,V\) 为有限维 \(k\) 向量空间，\(b:U\times V\to k\) 为双线性配对。令 \(L_b:U\to V^*\)、\(u\mapsto b(u,-)\) 及 \(R_b:V\to U^*\)、\(v\mapsto b(-,v)\)，则 \(b\) 非退化当且仅当 \(L_b\) 与 \(R_b\) 都是同构；这也等价于其中任一个是同构。此时 \(\dim U=\dim V\)。
\end{proposition}
\begin{proof}
非退化使两个映射都单射，故 \(\dim U\le\dim V\le\dim U\)，它们因而同为同构。若只知 \(L_b\) 是同构，且 \(R_b(v)=0\)，则全部 \(U\) 中元素都与 \(v\) 配为零。\(L_b\) 满射说明每个 \(V\) 上的线性泛函都在 \(v\) 处为零，由\cref{b2:lem:functional-extension}的分离结论，\(v=0\)。又两空间等维，所以 \(R_b\) 也是同构。交换两个变量证明另一方向。
\end{proof}

例如，求值 \(V^*\times V\to k\) 在有限维时完美。若取 \(b:k^2\times k^3\to k\)、\(b(u,v)=u_1v_1+u_2v_2\)，左根空间为零，右根空间却是一维；只检验一端不够。

\ProseParagraphBreak
对于同一空间上的一般型，条件 \(b(u,v)=0\) 与 \(b(v,u)=0\) 未必等价。以下两种型允许我们合并这些条件。

\begin{definition}\label{b2:def:symmetric-alternating-form}
设 \(k\) 为域，\(V\) 为 \(k\) 向量空间，\(b:V\times V\to k\) 为双线性型。如果 \(b(u,v)=b(v,u)\) 对所有 \(u,v\in V\) 成立，就称 \(b\) 为\term[duichen-xing]{对称型}[symmetric bilinear form]；如果 \(b(v,v)=0\) 对所有 \(v\in V\) 成立，就称它为\term[jiaoti-xing]{交替型}[alternating bilinear form]。如果 \(b(u,v)=-b(v,u)\) 对所有 \(u,v\in V\) 成立，就称它为\term[fanduichen-xing]{反对称型}[skew-symmetric bilinear form]。
\end{definition}

将 \(b(u+v,u+v)=0\) 展开，便知交替型总是反对称的。对于子空间 \(W\)，先分别记
\[
W^{\perp_r}=\{\,v\in V\mid b(w,v)=0\ (w\in W)\,\},\qquad
{}^{\perp_l}W=\{\,v\in V\mid b(v,w)=0\ (w\in W)\,\}.
\]
若 \(b\) 对称，或满足 \(b(v,u)=-b(u,v)\)，两者相同，统记为 \(W^\perp\)，称为\term[zhengjiao-bu]{正交补}[orthogonal complement]；此时 \(\operatorname{rad}b=V^\perp\) 称为型的\term[xing-de-gen]{根空间}[radical of a form]。这里的“补”首先是正交条件的名称，并不预先断言 \(V=W\oplus W^\perp\)。
\symidx{\(W^\perp\)：对称型或交替型下的正交补}
\symidx{\(\operatorname{rad}b\)：双线性型的根空间}

\begin{theorem}\label{b2:thm:orthogonal-complement}
设 \(k\) 为域，\(V\) 为有限维 \(k\) 向量空间，\(b:V\times V\to k\) 为非退化对称型或交替型，\(W\le V\)。记 \(W^\perp=\{v\in V\mid b(v,w)=0\text{ 对所有 }w\in W\}\)，则
\[
\dim W^\perp=\dim V-\dim W,\qquad (W^\perp)^\perp=W.
\]
若 \(b|_{W\times W}\) 非退化，则
\[
V=W\mathbin{\perp}W^\perp,
\]
其中符号 \(\perp\) 表示两个分量互相正交的直和；\(W^\perp\) 上的限制型也非退化。即使 \(b\) 在 \(V\) 上退化，只要它是对称型或交替型且 \(b|_{W\times W}\) 非退化，所述直和分解仍成立。
\end{theorem}
\begin{proof}
当 \(V\) 上非退化时，映射 \(V\to W^*\)、\(v\mapsto b(-,v)|_W\) 满射：先把 \(W\) 上的泛函延拓到 \(V\)，再用\cref{b2:prop:finite-perfect-pairing}。其核为 \(W^\perp\)，秩与核的维数公式给出第一个等式。正交条件的对称性给出 \(W\subseteq(W^\perp)^\perp\)，再由维数相等得到相等。

现在只假设 \(W\) 上非退化。对任意 \(v\in V\)，\cref{b2:prop:finite-perfect-pairing}在 \(W\) 上给出唯一 \(w\in W\)，使 \(b(x,w)=b(x,v)\) 对所有 \(x\in W\) 成立。因此 \(v-w\in W^\perp\)，而 \(W\cap W^\perp=0\) 又正是限制型非退化的条件。这证明直和分解。若 \(V\) 上也非退化，而 \(z\in W^\perp\) 与全部 \(W^\perp\) 正交，则它还与 \(W\) 正交，故与 \(V\) 正交，只能为零。
\end{proof}

例如，实平面上 \(b(x,y)=x_1y_1-x_2y_2\) 非退化，但直线 \(W=\mathbb R(1,1)\) 满足 \(W^\perp=W\)。维数公式正确，直和分解却不成立，因为限制型为零。非退化性不能从整个空间自动传给每个子空间。

\subsection{对称型与交替型}
\label{b2:subsec:1001:02}

交替性已经蕴含反对称性。特征不为二时，反对称性反过来给出 \(2b(v,v)=0\)，所以二者等价。特征二时不能这样约去二；对称型 \(b(x,y)=xy\) 就不是交替型。交替型的定义在任意特征下都要求对角取值为零。

\begin{theorem}[交替型的标准形]\label{b2:thm:alternating-form-normal}
设 \(k\) 为任意特征的域，\(V\) 为有限维 \(k\) 向量空间，\(b:V\times V\to k\) 为交替双线性型，则 \(V\) 有一组基
\[
e_1,f_1,\ldots,e_m,f_m,z_1,\ldots,z_r,
\]
使不同的两维分量及最后的根空间互相正交，并有
\[
b(e_i,f_i)=1,\quad b(f_i,e_i)=-1,\quad b(e_i,e_i)=b(f_i,f_i)=0.
\]
\(z_1,\ldots,z_r\) 是根空间的基。因此型的秩为 \(2m\)，并完全决定其在给定维数下的合同类型。特别地，非退化交替型只存在于偶数维空间。
\end{theorem}
\begin{proof}
若型为零，任取基即可。否则选 \(e,f_0\) 使 \(b(e,f_0)\ne0\)，将 \(f_0\) 除以该值，得到 \(b(e,f)=1\)。\(e,f\) 线性无关：若 \(ae+cf=0\)，分别与 \(e,f\) 配对便得 \(c=a=0\)。它们的线性生成空间 \(H\) 上的矩阵为
\[
\begin{pmatrix}0&1\\-1&0\end{pmatrix},
\]
行列式为一，故限制型非退化。由\cref{b2:thm:orthogonal-complement}中不要求整体非退化的分解，\(V=H\perp H^\perp\)。在较低维的 \(H^\perp\) 上归纳，直到剩余限制型为零。每个已分离的两维块非退化，故最后的零型分量恰是整个型的根空间。

上述基给出的矩阵由 \(m\) 个同样的可逆两维块及一个 \(r\) 维零块组成，所以秩为 \(2m\)。秩不随换基改变，\(r=\dim V-2m\) 也被确定；同秩的两个型把上述基逐一对应便得到保持配对的同构。
\end{proof}

非退化交替型的一组 \(e_i,f_i\) 称为\term[xin-ji]{辛基}[symplectic basis]。若先排列 \(e_1,\ldots,e_m\)，再排列 \(f_1,\ldots,f_m\)，其矩阵就是
\[
J_m=\begin{pmatrix}0&I_m\\-I_m&0\end{pmatrix}.
\]
这些论证没有除以二，因而在特征二中仍然有效；此时 \(-I_m=I_m\)，但对角线仍为零。

\subsection{合同分类与相似分类}
\label{b2:subsec:1001:03}

选基 \(\mathcal E=(e_1,\ldots,e_n)\)，令 \(B=\Bigl(b(e_i,e_j)\Bigr)\)。若 \(x,y\) 是两个向量的坐标列，则
\[
b(u,v)=x^{\mathsf T}By.
\]
这既说明矩阵由型唯一确定，也说明每个矩阵都确定一个双线性型。非退化恰好等于 \(B\) 可逆；对称性等于 \(B^{\mathsf T}=B\)，交替性等于 \(B^{\mathsf T}=-B\) 且对角元全零。

\begin{proposition}\label{b2:prop:bilinear-congruence}
设 \(k\) 为域，\(V\) 为 \(n\) 维 \(k\) 向量空间，\(b:V\times V\to k\) 为双线性型，其旧基矩阵为 \(B\)。若新基的列在旧基中的坐标组成 \(P\in\operatorname{GL}_n(k)\)，则型的新矩阵为
\[
B'=P^{\mathsf T}BP.
\]
存在这样的 \(P\) 时，称两矩阵\term[hetong]{合同}[congruent]。两个型之间的\term[dengju-tonggou]{等距同构}[isometry]是保持双线性取值的线性双射；它们等距，当且仅当相应矩阵合同。
\end{proposition}
\begin{proof}
新坐标为 \(x',y'\) 时，旧坐标为 \(Px',Py'\)，故取值是 \((x')^{\mathsf T}P^{\mathsf T}BPy'\)。逐个代入标准基向量便得到矩阵公式。一般线性双射也由可逆矩阵给出，保持两变量取值的等式正是同一个合同等式。
\end{proof}

算子换基的公式则为 \(A'=P^{-1}AP\)。算子的来源与目标是同一个空间，型的两个位置都是接受向量的变量，因此一个公式出现逆矩阵，另一个出现转置。实一维型的矩阵 \((1)\) 与 \((4)\) 通过 \(P=(2)\) 合同，但一维矩阵相似只能使其不变，故二者不相似。

\ProseParagraphBreak
合同保持秩。对于非退化型，等式
\[
\det B'=(\det P)^2\det B
\]
还说明行列式在 \(k^\times/(k^\times)^2\) 中的平方类不变。不过一般域上的平方类不是完整分类数据；例如实数域上 \(I_2\) 与 \(-I_2\) 行列式相同，却一个恒取正值，一个恒取负值。下一节将从二次取值出发说明这里还缺少什么信息。
